\section{Question 1}

\subsection{Question}
Demonstrate that you know how to use "curl" well enough to 
correctly POST data to a form.  Show that the HTML response that
is returned is "correct".  That is, the server should take the
arguments you POSTed and build a response accordingly.  Save the
HTML response to a file and then view that file in a browser and
take a screen shot.

\subsection{Resources}
\begin{itemize}
\item \LaTeX: \url{http://www.electronics.oulu.fi/latex/examples/example_1}
\item \LaTeX: \url{http://scott.sherrillmix.com/blog/programmer/displaying-code-in-latex/}
\item curl: \url{http://curl.haxx.se/docs/httpscripting.html#POST}
\end{itemize}

\subsection{Answer}
Using a simple wiki server built from the Go language net/http package tutorial, found here: \url{https://golang.org/doc/articles/wiki/}
\begin{lstlisting}
[mchaney@mchaney-d gowiki]$ curl -v -d "body=something something" localhost:8080/save/TestPage
* About to connect() to localhost port 8080 (#0)
*   Trying ::1...
* connected
* Connected to localhost (::1) port 8080 (#0)
> POST /save/TestPage HTTP/1.1
> User-Agent: curl/7.27.0
> Host: localhost:8080
> Accept: */*
> Content-Length: 24
> Content-Type: application/x-www-form-urlencoded
> 
* upload completely sent off: 24 out of 24 bytes
< HTTP/1.1 302 Found
< Location: /view/TestPage
< Date: Fri, 29 Aug 2014 13:32:01 GMT
< Content-Length: 0
< Content-Type: text/plain; charset=utf-8
< 
* Connection #0 to host localhost left intact
* Closing connection #0
\end{lstlisting}

The wiki uses the {\tt/save/} url to save the contents of a text box as the body of the target page. 

\begin{figure}[h!]
\centering
\label{fig:curl_ss}
\fbox{\includegraphics[scale=.45]{q1/curl.png}}
\caption{Screenshot of curl POST result}
\end{figure}